%===========================================================================
% CONFIGURACIÓN DEL DOCUMENTO
%===========================================================================
\documentclass[12pt,a4paper,spanish,twoside]{article}
\usepackage[a4paper,top=3cm,bottom=3cm,left=3cm,right=3cm]{geometry}

% Separación de los párrafos
\setlength{\parskip}{0.6cm}

\usepackage[utf8]{inputenc}
\usepackage[spanish,activeacute]{babel}
\usepackage[T1]{fontenc}
\usepackage{wrapfig}
\usepackage{url}
\usepackage{epsfig}
\usepackage{titlesec}

% ─────────────────────────────────────────────
% PAQUETE siunitx — configuración para español
% ─────────────────────────────────────────────
\usepackage{siunitx}
\sisetup{
    group-separator   = {.},         % separador de miles: punto
    group-minimum-digits = 4,        % agrupar desde 4 dígitos (1000 → 1.000)
    output-decimal-marker = {,},     % por si locale no lo aplica solo
    per-mode          = symbol,      % a/b en lugar de a b^{-1}
}
\DeclareSIUnit{\bit}{bit}
\DeclareSIUnit{\byte}{byte}
\DeclareSIUnit{\GiB}{GiB}

% --- Paquetes añadidos para el TFM (no forman parte de la plantilla original) ---
\usepackage{booktabs}         % Líneas de tabla profesionales (\toprule, \midrule, \bottomrule)
\usepackage{tabularx}         % Columnas de ancho variable en tablas
\usepackage{amsmath}          % Entornos matemáticos extendidos
\usepackage{amssymb}          % Símbolos matemáticos adicionales (\mathbb, \mathcal, etc.)
\usepackage{graphicx}         % Soporte para imágenes PNG, JPG, PDF
\usepackage{xcolor}           % Necesario para los colores
\usepackage[numbers]{natbib}  % Para insertar citas narrativas,
\usepackage{placeins}         % Para controlar la posición de las figuras con \FloatBarrier

% Definición de colores
\definecolor{altarelevancia}{RGB}{46,125,50}    % verde oscuro
\definecolor{normalrelevancia}{RGB}{255,179,0}  % ámbar
\definecolor{marginalrelevancia}{RGB}{230,81,0} % naranja quemado
\definecolor{descartes}{RGB}{198,40,40}         % rojo
% ---------------------------------------------------------------------------------

% Tablas muy largas que ocupan más de una página (se parten)
\usepackage{longtable}

% Paquete para subrayado, tachado, etc.
\usepackage{ulem}
\normalem

% Incluir paragraph como nivel de sección numerado
\setcounter{secnumdepth}{4}
\setcounter{tocdepth}{4}
\makeatletter
\renewcommand\paragraph{\@startsection{paragraph}{4}{\z@}%
                                     {-3.25ex\@plus -1ex \@minus -.2ex}%
                                     {1.5ex \@plus .2ex}%
                                     {\normalfont\normalsize\bfseries}}
\makeatother

\usepackage{listings}
\usepackage{color}
\definecolor{listinggray}{gray}{0.9}
\lstset{language=Matlab}
\lstset{keywordstyle=\color{red}\bfseries\underbar}

% Encabezados y pie de página
\usepackage{fancyhdr}
\usepackage{fancyvrb}

% Incluir índices y bibliografía en el índice general
\usepackage[nottoc]{tocbibind}

\makeatother
\usepackage[pageref]{backref}
\renewcommand{\backreftwosep}{ y }
\renewcommand{\backreflastsep}{ y }

\renewcommand*{\backref}[1]{}
\renewcommand{\backrefalt}[4]{%
\ifcase #1 %
{\unskip\nobreak\hfil\penalty50\null\hfil%
\mbox{\nobreakspace\small[\textsl{No citado: usado como
referencia}]}%
\parfillskip=0pt\finalhyphendemerits=0\par}
\or
{\unskip\nobreak\hfil\penalty50\null\hfil%
\mbox{\nobreakspace\small[\textsl{Citado en
p\'ag.\nobreakspace#2.}]}%
\parfillskip=0pt\finalhyphendemerits=0\par}
\else
{\unskip\nobreak\hfil\penalty50\null\hfil%
\mbox{\nobreakspace\small[\textsl{Citado en
p\'ags.\nobreakspace#2.}]}%
\parfillskip=0pt\finalhyphendemerits=0\par}
\fi
}

\usepackage{setspace}
\setlength\headheight{14.6pt}


%===========================================================================
% INICIO DEL DOCUMENTO
%===========================================================================
\begin{document}

\renewcommand{\tablename}{Tabla}
\renewcommand{\listtablename}{Índice de tablas}

%===========================================================================
% PORTADA
%===========================================================================
\begin{titlepage}

\vspace{4.3cm}

\begin{flushright}
	{\large \bf Informe Técnico - Technical Report}\\
	{\large \bf DPTOIA-IT-BORRADOR}\\
	{\large \bf \today}
\end{flushright}

\begin{center}
\vspace{5.2cm}
%-----------------------------------------------------------
% CONTENIDO: Título del informe técnico
{\Large \bf Título del Informe Técnico}\\
%-----------------------------------------------------------

\vspace{2cm}
%-----------------------------------------------------------
% CONTENIDO: Nombres de los autores
{\bf Autor 1}\\
{\bf Autor 2}\\
{\bf Autor 3}\\
{\bf Autor 4}\\
%-----------------------------------------------------------
\end{center}

\vspace{3.9cm}
\begin{figure}[h!]
    \begin{center}
    \footnotesize
        \includegraphics[width=3cm]{Imagen1-eps-converted-to.pdf}
        \includegraphics[width=3cm]{Imagen2-eps-converted-to.pdf}
    \end{center}
\end{figure}

\vspace{1.5cm}
    \begin{center}
        {\bf Departamento de Informática y Automática}\\
        {\bf Universidad de Salamanca}
    \end{center}
\end{titlepage}


%===========================================================================
% PÁGINA DE REVISORES E INFORMACIÓN DE AUTORES 
%===========================================================================
\thispagestyle{empty}
\vspace{4.9cm}
\begin{flushleft}
Revisado por:\\
\hspace{1cm}{Nombres de los autores}\\
\hspace{1cm}{Nombres de los autores}

\vspace{1cm}
Información del Autor:\\[0.5cm]

%-----------------------------------------------------------
% CONTENIDO: Datos de los autores (nombre, afiliación, correo)
\hspace{1cm}{\footnotesize Autor 1}\\
\hspace{1.5cm}{\footnotesize 1ª Línea Autor1}\\
\hspace{1.5cm}{\footnotesize 2ª Línea Autor1}\\
\hspace{1.5cm}{\footnotesize 3ª Línea Autor1}\\
\hspace{1.5cm}{\footnotesize correo@autor1.com}\\[0.5cm]

\hspace{1cm}{\footnotesize Autor 2}\\
\hspace{1.5cm}{\footnotesize 1ª Línea Autor2}\\
\hspace{1.5cm}{\footnotesize 2ª Línea Autor2}\\
\hspace{1.5cm}{\footnotesize 3ª Línea Autor2}\\
\hspace{1.5cm}{\footnotesize correo@autor2.com}\\[0.5cm]

\hspace{1cm}{\footnotesize Autor 3}\\
\hspace{1.5cm}{\footnotesize 1ª Línea Autor3}\\
\hspace{1.5cm}{\footnotesize 2ª Línea Autor3}\\
\hspace{1.5cm}{\footnotesize 3ª Línea Autor3}\\
\hspace{1.5cm}{\footnotesize correo@autor3.com}\\[0.5cm]

\hspace{1cm}{\footnotesize Autor 4}\\
\hspace{1.5cm}{\footnotesize 1ª Línea Autor4}\\
\hspace{1.5cm}{\footnotesize 2ª Línea Autor4}\\
\hspace{1.5cm}{\footnotesize 3ª Línea Autor4}\\
\hspace{1.5cm}{\footnotesize correo@autor4.com}\\[0.5cm]
%-----------------------------------------------------------

\vfill
Este documento puede ser libremente distribuido.\\
(c) 2006 Departamento de Informática y Automática - Universidad de Salamanca.
\end{flushleft}

%===========================================================================
% DECLARACIÓN DE AUTORÍA
%===========================================================================«»
\newpage
\thispagestyle{empty}

\vspace*{0.5cm}
    \begin{center}
        {\textbf{Declaración de Autoría}}
    \end{center}
\vspace{0.5cm}

\noindent Declaro que he redactado \underline{\hspace{1.5cm}} titulado \underline{\hspace{1.5cm}} del \underline{\hspace{1.5cm}} en el \underline{\hspace{1.5cm}} del curso académico \underline{\hspace{1.5cm}} de forma autónoma, con la ayuda de las fuentes y la literatura citadas en la bibliografía, y que he identificado como tales todas las partes tomadas de las fuentes y de la literatura indicada, textualmente o conforme a su sentido.

\vspace{0.5cm}

\noindent Además, soy conocedor de que \underline{\hspace{1.5cm}} forma parte de los trabajos de investigación que llevan a cabo mis directores \underline{\hspace{1.5cm}} dentro de los grupos de investigación \underline{\hspace{1.5cm}} de la Universidad de Salamanca y, en consecuencia, comparto con ellos la propiedad intelectual de los resultados alcanzados.

\vspace{0.5cm}

\noindent En Salamanca, a \underline{\hspace{1.5cm}} de \underline{\hspace{4cm}} de \underline{\hspace{1.5cm}}

\vspace{1cm}

\noindent Fdo.: \rule{8cm}{0.4pt}\\[0.3cm]
\hspace*{1.5cm} \texttt{Autor}

\vfill
\begin{center}
    \footnotesize
    Facultad de Ciencias. Plaza de los Caídos, s/n, 37008 Salamanca.\\
    Tel.: +34 923 29 44 00 ext.: 1513 $\cdot$ Fax: +34 923 29 45 14\\
    \texttt{http://mastersi.usal.es} $\cdot$ \texttt{mastersi@usal.es}
\end{center}

%===========================================================================
% AGRADECIMIENTOS
%===========================================================================
\newpage
\thispagestyle{empty}

\vspace*{0.5cm}
    \begin{center}
        {\bf Agradecimientos}
    \end{center}
\vspace{0.5cm}

%-----------------------------------------------------------
% CONTENIDO: Texto de agradecimientos
En primer lugar, quisiera expresar mi más sincero agradecimiento ...

%-----------------------------------------------------------

\vfill
\begin{flushright}
    \textit{Firma del Autor}\\
    \textit{Lugar, \the\year}
\end{flushright}


%===========================================================================
% RESUMEN / ABSTRACT 
%===========================================================================
\newpage
\renewcommand{\thepage}{\roman{page}}
\setcounter{page}{1}
\pagestyle{fancy}
\fancyhf{}
\fancyfoot[LE,RO]{\textbf{\thepage}}
\fancyfoot[LO,RE]{\textbf{DPTOIA-IT-BORRADOR}}
\renewcommand{\headrulewidth}{0.0pt}
\renewcommand{\footrulewidth}{0.4pt}

\vspace{6.5cm}
\textbf{Resumen}

%-----------------------------------------------------------
% CONTENIDO: Resumen en castellano (150-250 palabras)
Se incluye el resumen en castellano.
%-----------------------------------------------------------

\vspace{1cm}
\textbf{Abstract}

%-----------------------------------------------------------
% CONTENIDO: Resumen en inglés (150-250 palabras)
Se incluye el resumen en inglés.
%-----------------------------------------------------------


%===========================================================================
% ÍNDICES
%===========================================================================
\newpage
\setlength{\parskip}{0.8mm}
\tableofcontents
\newpage
\listoffigures
\newpage
\listoftables
\setlength{\parskip}{2.5mm}

\newpage


%===========================================================================
% CUERPO DEL DOCUMENTO
%===========================================================================
\renewcommand{\thepage}{\arabic{page}}
\ifodd\thepage
\else
  .
  \thispagestyle{empty}
  \newpage
\fi
\setcounter{page}{1}
\fancyhf{}
\renewcommand{\headrulewidth}{0.4pt}

%-----------------------------------------------------------
% CONTENIDO: Sustituir "Título corto" y "Apellido et al."
\fancyhead[LE]{\textbf{Título corto del informe}}
\fancyhead[RO]{\textbf{Apellido et al.}}
%-----------------------------------------------------------

\fancyfoot[LE,RO]{\textbf{\thepage}}
\fancyfoot[LO,RE]{\textbf{DPTOIA-IT-BORRADOR}}


%===========================================================================
% INTRODUCCIÓN
%===========================================================================
\section{Introducción}
\label{sec:introduccion}

Texto intro, \cite{kreher05pseudocode}. O también, \citet{ohl97emp}.

\newpage

%===========================================================================
% BIBLIOGRAFÍA
%===========================================================================
\bibliographystyle{plainnat}
\bibliography{Bibliografia}

%===========================================================================
% ANEXO A
%===========================================================================
\newpage
\appendix
\section{Anexo A}
\label{anexo:Anexo}

\end{document}
